% Generated draft: requires mathematical and source review.
\documentclass[11pt]{article}
\usepackage[margin=25mm]{geometry}
\usepackage{fontspec}
\setmainfont{Latin Modern Roman}
\usepackage{amsmath,amssymb,mathtools,unicode-math}
\setmathfont{Latin Modern Math}
\usepackage{enumitem,needspace,xurl,hyperref}
\setlength{\parindent}{0pt}
\setlength{\parskip}{5pt}
\setlength{\emergencystretch}{3em}
\newcommand{\N}{\mathbb{N}}
\newcommand{\Z}{\mathbb{Z}}
\newcommand{\Q}{\mathbb{Q}}
\newcommand{\R}{\mathbb{R}}
\newcommand{\C}{\mathbb{C}}
\newcommand{\F}{\mathbb{F}}
\newcommand{\A}{\mathbb{A}}
\newcommand{\PP}{\mathbb{P}}
\newcommand{\E}{\mathbb{E}}
\newcommand{\eps}{\varepsilon}
\newcommand{\dd}{\,\mathrm{d}}
\begin{document}
\section*{UnsolvedMath ID: 30006476}
\textbf{Source problem code:} catalogue entry
\subsection{Definitions and mathematical statement}
Let \(K\) be a number field and \(G\) a finite solvable group. Interpret \(G\) as the constant finite group scheme over \(K\); this convention makes explicit the usual interpretation of the supplied statement, which specifies no Galois action. Write \(H^1(K,G)\) for the pointed set of isomorphism classes of \(G\)-torsors over \(\operatorname{Spec}K\). Equivalently, if \(K^{\mathrm{sep}}\) is a separable closure of \(K\), then
\[
H^1(K,G)=\operatorname{Hom}_{\mathrm{cont}}\!\left(\operatorname{Gal}(K^{\mathrm{sep}}/K),G\right)/G,
\]
where \(G\) acts by conjugation; surjectivity of the homomorphisms is not required.

Use the following standard convention for \(R\)-equivalence. Two classes \(\alpha,\beta\in H^1(K,G)\) are directly \(R\)-equivalent if there are a Zariski open subset \(U\subseteq\mathbb A^1_K\) containing \(0\) and \(1\), and a \(G\)-torsor \(\mathcal T\to U\), whose fibers at \(0\) and \(1\) represent \(\alpha\) and \(\beta\), respectively. Let \(R\) be the equivalence relation generated by direct \(R\)-equivalence, and set
\[
r(K,G):=\#\bigl(H^1(K,G)/R\bigr).
\]
The supplied extract does not itself define \(R\)-equivalence, so this convention requires confirmation against the source.

\textbf{Finiteness question.} Is it true that, for every number field \(K\) and every finite solvable group \(G\),
\[
r(K,G)<\infty?
\]

\textbf{Additional requested relation, unresolved in the supplied evidence.} If \(r(K,G)\) is finite, the supplied statement asks for its relation to the minimum number of \(G\)-torsors over a purely transcendental extension \(K(t_1,\ldots,t_d)\) needed to densely parametrize \(H^1(K,G)\). Neither the meaning of dense parametrization nor the allowed specializations, the quantification of \(d\), or the proposed relation is specified. Thus this additional request cannot be made into a precise assertion from the supplied record alone. The extract contains a preceding claim involving a group \(A\), but omits its hypotheses and ends with the incomplete phrase ``for all sufficiently large''; that claim cannot be transferred to \(G\).

The supplied literature assessment also mentions preservation of finiteness under split extensions. It provides no exact formulation or hypotheses for that additional question.
\subsection{Short English statement}
For every number field \(K\) and finite solvable constant group \(G\), ask whether there are finitely many classes of \(G\)-torsors under the equivalence relation generated by one-parameter families joining their fibers. The requested connection with dense parametrization, and the reference to split extensions, need further source clarification. The supplied open status is dated curation.
\subsection{Sources}
\begin{itemize}
\item Local UnsolvedMath registry, \texttt{data/unsolvedmath/problems.json}, permanent ID \(30006476\), source problem code \texttt{catalogue entry}, titled ``Finiteness of R-Equivalence Classes in Nonabelian Cohomology.'' This draft uses the supplied complete registry record; its supplied statement status is \texttt{corrected\_verified}, with repairs described as limited to notation, grammar, and typography.
\item Supplied source citation: \emph{Arithmetic Statistics for Algebraic Objects} (2026), Oberwolfach Reports, DOI 10.4171/owr/2025/52, pp. 2765--2842. Supplied source URL and literature reference ``Oberwolfach report 52/2025'': \url{https://doi.org/10.4171/owr/2025/52}. The extract identifies the question as Problem 5. No proposer is supplied.
\item Literature assessment dated exactly \texttt{2026-08-22}: the supplied curation labels the question open and describes finiteness for finite solvable groups over number fields, including preservation under split extensions, as explicitly posed. This is supplied dated curation, not an independently verified assessment of the literature or current status.
\end{itemize}
\end{document}
